\documentclass[xcolor=x11names,compress]{beamer}

%% General document %%%%%%%%%%%%%%%%%%%%%%%%%%%%%%%%%%
\usepackage{graphicx}
\graphicspath{{graphics/}}
\usepackage{animate}
\usepackage{tikz}

% ########For highlight boxes over equations########
\usepackage{color}
\input{inputs/rgb}
\definecolor{capri}{rgb}{0,.75,1}
\definecolor{coral}{rgb}{1,.2,.2}
\newcommand{\hltyellow}[1]{\colorbox{yellow}{$\displaystyle #1$}}
\newcommand{\hltblue}[1]{\colorbox{capri}{$\displaystyle #1$}}
\newcommand{\hltred}[1]{\colorbox{coral}{$\displaystyle #1$}}
\newcommand{\hltgreen}[1]{\colorbox{green}{$\displaystyle #1$}}

\usepackage{ulem}
\renewcommand<>{\sout}[1]{
  \only#2{\beameroriginal{\sout}{#1}}
  \invisible#2{#1}
}

%%%%%%%%%%%%%%%%%%%%%%%%%%%%%%%%%%%%%%%%%%%%%%%%%%%%%%

%% Beamer Layout %%%%%%%%%%%%%%%%%%%%%%%%%%%%%%%%%%

\usetheme{Madrid}
\usecolortheme{seahorse}
\setbeamertemplate{headline}{}

% ###Comment/Uncomment to display/hide navigation symbols at bottom:
\setbeamertemplate{navigation symbols}{}

% ###Comment/Uncomment to hide/display outline at top:
% \useoutertheme[subsection=false,shadow]{miniframes}

% \useinnertheme{default}

\setbeamertemplate{footline}
{
  \leavevmode%
  \hbox{
  \begin{beamercolorbox}[wd=.1\paperwidth,ht=2.2ex,dp=1ex,center]{author in head/foot}
    \usebeamerfont{author in head/foot}\insertshortauthor
  \end{beamercolorbox}
  \begin{beamercolorbox}[wd=.8\paperwidth,ht=2.2ex,dp=1ex,center]{title in head/foot}%
    \usebeamerfont{title in head/foot}\insertshorttitle\hspace*{3em}
  \end{beamercolorbox}}
  \begin{beamercolorbox}[wd=.05\paperwidth,ht=2.2ex,dp=1ex,center]{}
     \insertframenumber{} / \inserttotalframenumber\hspace*{-3ex}
  \end{beamercolorbox}
}

\title[VectorBiTE Training: NLLS]{\bf Fitting mathematical models to data using Non-linear Least-Squares}
% \subtitle{Silwood Masters Course Lecture}
\author[Samraat]{Samraat Pawar}
\institute{\it
	{Department of Life Sciences\\
	\vspace{5pt}
	(Silwood Park)\\
		% \vspace{5pt}
% Imperial College London
	}
	\vspace{20pt}
  \centering
%   \includegraphics[height = .3in]{graphics/Imperial_Color1.pdf}
  \includegraphics[height = .5in]{../../notebooks/graphics/VB_logo.jpg}

}

\usefonttheme{serif}
\usepackage{palatino}

\setbeamerfont{title like}{shape=\scshape}
\setbeamerfont{frametitle}{shape=\scshape, series = \bfseries}
\setbeamertemplate{frametitle}[default][center]

\setbeamercolor*{lower separation line head}{bg=DeepSkyBlue4} 
\setbeamercolor*{normal text}{fg=black,bg=white} 
\setbeamercolor*{alerted text}{fg=red} 
\setbeamercolor*{example text}{fg=black} 
\setbeamercolor*{structure}{fg=black}
 
\setbeamercolor*{palette tertiary}{fg=black,bg=black!10} 
\setbeamercolor*{palette quaternary}{fg=black,bg=black!10} 

\renewcommand{\(}{\begin{columns}}
\renewcommand{\)}{\end{columns}}
\newcommand{\<}[1]{\begin{column}{#1}}
\renewcommand{\>}{\end{column}}
  
%%%%%%%%%%%%%%%%%%%%%%%%%%%%%%%%%%%%%%%%%%%%%%%%%%%%
%%%%%%%%%%%%%%%%%%%%%%%%%%%%%%%%%%%%%%%%%%%%%%%%%%%%

\begin{document}

%%%%%%%%%%%%%%%%%%%%%%%%%%%%%%%%%%%%%%%%%%%%%%%%%%%%%%
\begin{frame}[plain]
 
\titlepage
\end{frame}

%%%%%%%%%%%%%%%%%%%%%%%%%%%%%%%%%%%%%%%%%%%%%%%%%%%%%%
\begin{frame}{Outline}

\begin{itemize}\itemsep16pt
	\item Why Non-Linear Least Squares regression / fitting? 
	\item The NLLS fitting method 
	\item NLLS in {\tt R}
	\item Afternoon practicals overview
\end{itemize}

\end{frame}

%%%%%%%%%%%%%%%%%%%%%%%%%%%%%%%%%%%%%%%%%%%%%%%%%%%%%%
\begin{frame}{Linear models are great}

\includegraphics[width=\textwidth]{Linear.pdf}

\begin{itemize}[<+->]\itemsep4pt
	\pause 
	\item These are {\it all} good {\it Linear Models} (really?!)
	\item The data can be modelled (aka "fitted to a mathematical model") as a {\it linear combination} of {\it variables} and {\it coefficients}
	\item Easily fitted using {\it Ordinary Least Squares} (OLS) regression
	\item Linear models can {\it include curved responses} (e.g. polynomial regression)
	\item OK, so then {\it why Non-Linear Least Squares (NLLS) fitting?} 
\end{itemize}

\end{frame}

% %%%%%%%%%%%%%%%%%%%%%%%%%%%%%%%%%%%%%%%%%%%%%%%%%%%%%%
\begin{frame}{Why NLLS? -- first, what makes a model non-linear?}

\begin{itemize}[<+->]\itemsep10pt

	\item OLS can be used to fit both linear and nonlinear {\it equations} that {\it intrinsically linear}, e.g.,
	\begin{itemize}
		\item 	Straight line: $y_i = \beta_0 + \beta_1 x_i + \varepsilon_i$
		\item Polynomial: $y_i = \exp(\beta_0) + \beta_1 x_i + \beta_2 x_i^2 
		+ \varepsilon_i$
	\end{itemize} 	
	\item Indeed, for OLS to work, we need {\it intrinsic linearity} --- 
	i.e., the equation to be fitted (model) should be {\it linear in the parameters}	
	\item Are these models linear in their {\it parameters}?  
	\begin{itemize}
		\item $y_i = \beta_0 + \beta_1 x_i^{\beta_2} + \varepsilon_i$
		\item $y_i = \beta_0 e^{\beta_2 x_i} + \varepsilon_i$
	\end{itemize}
	\pause NO!
		
\end{itemize}

\end{frame}

%%%%%%%%%%%%%%%%%%%%%%%%%%%%%%%%%%%%%%%%%%%%%%%%%%%%%%
\begin{frame}{So what --- why is intrinsic non-linearity a problem?}

{\bf Recall what the Least Squares method does:}

	\begin{itemize}[<+->] \itemsep10pt

		\item Consider a predictor $x$, data $y$, $n$ observations, and a 
		model that we want to fit to the data: \\
		
		$f(x_i,\boldsymbol\beta) + \varepsilon_i$ \\ 
		
		where $\boldsymbol\beta = (\beta_1, \beta_2, \ldots,\beta_k)$ are 
		the model's $k$ parameters
	
		\item The objective is to find estimates of values of the $k$ 
		parameters ($\hat{\beta_j}$) that minimize the sum ($S$) of squared 
		residuals ($r_i$) (AKA RSS):
		
		$S = \sum\nolimits_{i=1}^n{\left[ y_i - f(x_i,\boldsymbol\beta) 
		\right]^2} = \sum\nolimits_{i=1}^n{r_i^2}$
	
	\end{itemize}

\end{frame}

%%%%%%%%%%%%%%%%%%%%%%%%%%%%%%%%%%%%%%%%%%%%%%%%%%%%%%
\begin{frame}{The Least-squares solution}

	OLS minimizes the {\it sum} of the {\it squared} residuals
	
	\begin{columns}[T]
	\column{\textwidth}
		\animategraphics[width=\textwidth, controls, buttonsize=1em, %timeline=Varying_timeline.txt % does not work - TODO
		]{10}{Varying}{}{}
	\end{columns}
	
	\end{frame}

	
%%%%%%%%%%%%%%%%%%%%%%%%%%%%%%%%%%%%%%%%%%%%%%%%%%%%%%
\begin{frame}{If the model is linear, the solution is easy using algebra}
	
\begin{columns}[T]

	\column{0.5\textwidth}
		\includegraphics[width=\textwidth]{JustRight.pdf}
		
	\column{0.5\textwidth}
		\begin{align*}
		  y_1  & = \beta_0 + \beta_1 x_1  + {\color{red} \varepsilon_i }\\
		  \\
		  9.50  &= 5 + 4 \times 1 + {\color{red} 0.50}\\
		  11.00 &= 5 + 4 \times 2 - {\color{red} 2.00}\\
		  19.58 &= 5 + 4 \times 3 + {\color{red} 2.58}\\
		  20.00 &= 5 + 4 \times 4 - {\color{red} 1.00} 
		\end{align*}
		
		\[\beta_0 = 5; \beta_1 = 4\]
				
\end{columns}		

\end{frame}

%%%%%%%%%%%%%%%%%%%%%%%%%%%%%%%%%%%%%%%%%%%%%%%%%%%%%%
\begin{frame}{Intrinsic non-linearity does not allow a algebraic solution}

\begin{itemize}[<+->]\itemsep10pt

	\item So, then, in an {intrinsically non-linear model} such as $y_i = 
	\beta_0 e^{\beta_2 x_i} + \varepsilon_i$ the derivatives 
	$\frac{\partial r_i}{\partial \beta_j}$ are naughty
	
	\item That is, they are functions of both $x$ and the parameters 
	$\beta_j$, so the gradient equations do not have a solution like the 
	OLS case

	\item So the nice trick of solving  $\mathbf{Y} = \mathbf{X 
	\boldsymbol\beta  +  \boldsymbol\varepsilon}$ is impossible {\it 
	mathematically}
							
\end{itemize}

\end{frame}

%%%%%%%%%%%%%%%%%%%%%%%%%%%%%%%%%%%%%%%%%%%%%%%%%%%%%%
\begin{frame}{So --- Enter NLLS}

But we can use brute-force computation to find close-to-optimal least squares minimization!
	
\begin{itemize}[<+->]\itemsep10pt
	\item Choose initial values for the parameters we want to estimate ( $\beta_j$'s)
	
	\item Then, ``refine'' the parameters {\it iteratively} by calculating 
	$\frac{\partial r_i}{\partial \beta_j}$ {\it approximately} --- 
	this approximation is the {\it Jacobian} (the gradient), which is a matrix of 
	the $\frac{\partial r_i}{\partial \beta_j}$'s

	\item Whether a refinement has taken place in any step of the 
	iteration is determined by re-calculating the residuals at that 
	step

	\item Eventually, if it all goes well, we find a combination of 
	$\beta_j$'s that is {\it very close} to the desired solution 
	$\frac{\partial S}{\partial \beta_j} = 0, j = 0,1,2,\ldots,k$
			
\end{itemize}

\end{frame}

% %%%%%%%%%%%%%%%%%%%%%%%%%%%%%%%%%%%%%%%%%%%%%%%%%%%%%%
\begin{frame}[plain]{}

\begin{columns}[T]
	\column{\textwidth}
		\animategraphics[width=\textwidth, controls, buttonsize=1em, %timeline=Varying_timeline.txt % does not work - TODO
		]{10}{Varying}{}{}
	\end{columns}

\end{frame}

%%%%%%%%%%%%%%%%%%%%%%%%%%%%%%%%%%%%%%%%%%%%%%%%%%%%%%
\begin{frame}{OK, fine, why would {\it I} ever need NLLS?}

\begin{itemize}[<+->]\itemsep10pt
	\pause 
	\item Many observations in biology are just {\it not} well-fitted by a linear model

	\item That is, the underlying biological phenomena/phenomenon are not 
	well-described by a linear equation
	
	\item Examples:

	\begin{itemize}\itemsep2pt
	
		\item Logistic population growth
		
		\item Allometric growth

		\item Michaelis-Menten biochemical kinetics (two parameters 
		$V_{\max}$ and $K_m$: $v = \frac{V_{\max}[S]}{K_{m} + [S]}$
	
		\item Responses of metabolic rates to changing temperature 

		\item Consumer-Resource (e.g., predator-prey) functional responses 
		
		\item Time-series data (e.g., fitting a sinusoidal function)
			
	\end{itemize}

	\item \it Can you think of some examples? 

\end{itemize}

\end{frame}

%%%%%%%%%%%%%%%%%%%%%%%%%%%%%%%%%%%%%%%%%%%%%%%%%%%%%%
\begin{frame}{Example: Temperature and metabolism}
      
\begin{columns}[c]
  \column{2.1in}
    \pause
      \includegraphics[width=\textwidth]{graphics/Photobacterium.pdf}
  \pause
  \column{2.7in}
    $B = B_0 \boxed {e^{-\frac{E}{kT}}}f(T,T_{pk},E_D)$\\
    \vspace{10pt}
    \raggedright{$T$ = temperature (K)\\
     $k$ = Boltzmann constant (eV K$^{-1}$)}\\
     $E$ = Activation energy (eV)\\
     $T_{pk}$ = Temperature of peak performance\\
     $E_D$ = Deactivation energy (eV)\\
    {\small (J H van’t Hoff 1884, S Arrhenius 1889)}
\end{columns}

\begin{itemize}\setlength{\itemindent}{0em}
  \pause
  \item Surely there is more to thermal responses?
    \begin{itemize}\setlength{\itemindent}{-1em}
     \item Oxygen limitation
     \item Complexity of metabolic network
     \item Hormonal regulation
    \end{itemize}
    \pause
  \item \textit{What about {\it alternative models}?}\\

\end{itemize}
 
\end{frame}

%%%%%%%%%%%%%%%%%%%%%%%%%%%%%%%%%%%%%%%%%%%%%%%%%%%%%%
\begin{frame}{Example: Functional responses}
 
\begin{columns}[c]
  \column{2in}
  \begin{center}
			   $f(x_R) = \frac{{a x_R^{q + 1}}}{{1 + h a x_R^{q + 1}}}$
	\scriptsize (Holling, 1959)
	\end{center}
 \scriptsize
 $x_R$ = Resource density (Mass / Area or Volume)\\
 $a$ = Search rate (Area or Volume / Time )\\
 $h$ = Handling time \\
 $q$ = Shape parameter (dimensionless) 
 \column{2.3in}
      \includegraphics[width=\textwidth]{graphics/FR.pdf}
\end{columns}

\vspace{10pt}
 Note that: 
\begin{itemize}
\item NLLS fitting can yield $h < 0$, $q < 0$, or both 
\item $h < 0$ is biologically impossible but indicates an upward curving response
\item $q < 0$ is biologically unlikely as it indicates a decline in 
search rate with resource density (but is useful as a measure of 
deviation away from a type III response) 
\end{itemize}

\end{frame}

%%%%%%%%%%%%%%%%%%%%%%%%%%%%%%%%%%%%%%%%%%%%%%%%%%%%%%

\begin{frame}{NLLS fitting}

So the general procedure is:

\begin{enumerate}[<+->] \itemsep4pt
	\item Start with an initial value for each parameter in the model
	\item Generate the curve defined by the initial values
	\item Calculate the residual sum-of-squares (rss)
	\item Adjust the parameters to make the curve come closer to the data 
	points. This the tricky part --- more on this in the next slide
	\item Adjust the parameters again so that the curve comes even closer 
	to the points (rss decreases)
	\item Repeat 4--5
	\item Stop simulations when the adjustments make virtually no 
	difference to the rss
\end{enumerate}


\end{frame}

%%%%%%%%%%%%%%%%%%%%%%%%%%%%%%%%%%%%%%%%%%%%%%%%%%%%%%

\begin{frame}{NLLS fitting}

The {\it tricky part --- adjust parameters to make curve come closer to 
the data points} (step 4) has two main algorithms that you can choose between:

\begin{itemize}[<+->]\itemsep10pt

	\item The Gauss-Newton algorithm is the default in the {\tt nls} 
	package (part of the {\tt stats} base package) --- good in many 
	cases, but doesn't work very well if the model is mathematically 
	weird (the optimization landscape is difficult) and the starting 
	values for parameters are far-off-optimal
	
	\item The Levenberg-Marquardt (LM) switches switches between 
	Gauss-Newton and ``gradient descent'' and is more robust against 
	starting values that are far-off-optimal --- available in {\tt R} 
	through the the {\tt minpack.lm} package {\small \url 
	{http://cran.r-project.org/web/packages/minpack.lm}}
	
	% %% \item It's an R front-end to a Fortran LM implementation of the 
	% %% MINPACK package
	
	\item The command is {\tt nlsLM}

\end{itemize}   

\end{frame}

%%%%%%%%%%%%%%%%%%%%%%%%%%%%%%%%%%%%%%%%%%%%%%%%%%%%%%
\begin{frame}{NLLS fitting}

\begin{itemize}[<+->] \itemsep6pt
\item Once the algorithm as converged (hopefully -- but you may be surprised 
how well it usually works), you need to get the goodness of fit measures

\item First, of course, examine the fits visually

\item Also, report the best-fit results, including:

	\begin{itemize}
	
		\item Sums of deviations of the data points from the final model 
		fit (final RSS)
		\item $R^2$
		\item Estimated coefficients
		\item For each coefficient, standard error (can be used for CI's), t-statistic 
		and corresponding (two-sided) p-value 
	\end{itemize} 

\item The function {\tt summary.nls} will give you all these measures

\item Remember, the precise parameter values you obtain will depend in part on 
the initial values chosen and the convergence criteria 

\item You may also want to compare multiple models.  

\end{itemize}

\end{frame}

%%%%%%%%%%%%%%%%%%%%%%%%%%%%%%%%%%%%%%%%%%%%%%%%%%%%%%
\begin{frame}{NLLS Assumptions}

NLLS-regression has all the assumptions of OLS-regression:
	\begin{itemize}[<+->]\itemsep10pt
		
		\item No (in practice, minimal) measurement error in explanatory 
		variable ($x$-axis variable)
		
		\item Data have constant normal variance --- errors in the $y$-axis 
		are homogeneously distributed over the $x$-axis range
		
		\item The measurement/observation error distribution is Gaussian 
		--- for example, what would the error distribution of this 
		non-linear model be: $y_i = \beta_0 e^{\beta_2 x_i} + 
		\varepsilon_i$ 
		
		\item What if the errors are not normal? --- use Maximum Likelihood 
		or Bayesian methods instead.
	
	\end{itemize}
\end{frame}

%%%%%%%%%%%%%%%%%%%%%%%%%%%%%%%%%%%%%%%%%%%%%%%%%%%%%%

\begin{frame}{More NLLS tips}

\begin{itemize}[<+->]\itemsep10pt
	\item You can use mixed-effects modelling with NLLS in R; the package is {\tt nlme}
	{\small \url 
	{https://stat.ethz.ch/R-manual/R-devel/library/nlme/html/nlme.html}}
	(You are probably stuck with the Gauss-Newton algorithm with 
	{\tt nlme} though)
	
	\item You can also use Python --- look up {\tt lmfit} \url{https://lmfit.github.io/lmfit-py/index.html}. \\ {\tt python} seems to have a better Levenberg-Marqualdt implementation than {\tt R} 
	
\end{itemize}   

\end{frame}

%%%%%%%%%%%%%%%%%%%%%%%%%%%%%%%%%%%%%%%%%%%%%%%%%%%%%%

\begin{frame}{Readings and Resources}

\begin{itemize}\itemsep10pt
	
	\item Motulsky, Harvey, and Arthur Christopoulos. Fitting models to 
	biological data using linear and nonlinear regression: a practical 
	guide to curve fitting. OUP USA, 2004. 
	
	\item Johnson, J. B. \& Omland, K. S. 2004 Model selection in ecology 
and evolution. Trends Ecol. Evol. 19, 101--108. 

\end{itemize}   

\end{frame}

%%%%%%%%%%%%%%%%%%%%%%%%%%%%%%%%%%%%%%%%%%%%%%%%%%%%%%
\begin{frame}{So what --- why is intrinsic non-linearity a problem?}

{\bf That's all well and good, but we can use maths instead of 
brute-force computation!}

	\begin{itemize}[<+->] \itemsep10pt

		\item Our model is $f(x_i,\boldsymbol\beta) + \varepsilon_i$ \\ 
			
		\item We want to find estimates of values of the parameters 
		($\hat{\beta_j}$) that {\it minimize} the sum ($S$) of squared residuals 
		($r_i$) (AKA ``RSS'') $S = \sum\nolimits_{i=1}^n{\left[ y_i - 
		f(x_i,\boldsymbol\beta) \right]^2} = \sum\nolimits_{i=1}^n{r_i^2}$

		\item For this we can solve $\frac{\partial S}{\partial 
		\beta_j} = 0, j = 0,1,2,\ldots,k$ to find the {\it minimum}
	
		\item That is, we need to solve	
		$\frac{\partial \sum\nolimits_{i=1}^n{r_i^2}}{\partial \beta_j} = 
		0$
		
		\item Or, $2 \sum\nolimits_{i=1}^n r_i \frac{\partial r_i}{\partial \beta_j} = 0$
	
	\end{itemize}

\end{frame}

%%%%%%%%%%%%%%%%%%%%%%%%%%%%%%%%%%%%%%%%%%%%%%%%%%%%%%
\begin{frame}{So what --- why is intrinsic non-linearity a problem?}

	\begin{itemize}[<+->] \itemsep10pt

		\item Thus, solving $ 2 \sum\nolimits_{i=1}^n r_i \frac{\partial r_i}{\partial \beta_j} 
		= 0$ \\
		
		boils down to finding the ``gradient'' $ \frac{\partial r_i}{\partial \beta_j} $
	
		\item This is not a problem in linear models, because this gradient 
		is fully solvable as the equation is {\it intrinsically linear}
		
		\item That is, the solution of $ \frac{\partial r_i}{\partial 
		\beta_j} $ is simple (enough)  
				
	\end{itemize}

\end{frame}

%%%%%%%%%%%%%%%%%%%%%%%%%%%%%%%%%%%%%%%%%%%%%%%%%%%%%%
\begin{frame}{So what --- why is intrinsic non-linearity a problem?}

	\begin{itemize}[<+->] \itemsep15pt
		
		\item For example, if $f(x_i,\boldsymbol\beta) + \varepsilon_i = \beta_0 + \beta_1 x_i + 
		\varepsilon_i$ \\
		
		\item That is, our model is $y_i = \beta_0 + \beta_1 x_i + \varepsilon_i$ 
		(Linear Regression)
		
	\item Then we want to solve 
		
		$\frac{\partial S}{\partial {\beta_0}} = 
		\sum\nolimits_{i=1}^{n}{\frac{\partial {{\left[ 
		y_i-(\beta_0+\beta_1 x_i) \right]}^{2}}}{\partial {{\beta 
		}_{0}}}}=0 $ \\
		 
		$ \frac{\partial S}{\partial {{\beta 
		}_{1}}}=\sum\nolimits_{i=1}^{n}{\frac{\partial {{\left[ 
		{{y}_{i}}-({{\beta }_{0}}+{{\beta }_{1}}{{x}_{i}}) 
		\right]}^{2}}}{\partial {{\beta }_{1}}}} = 0$
		
	\end{itemize}

\end{frame}

%%%%%%%%%%%%%%%%%%%%%%%%%%%%%%%%%%%%%%%%%%%%%%%%%%%%%%
\begin{frame}{So what --- why is intrinsic non-linearity a problem?}

	\begin{itemize}[<+->] \itemsep12pt

		\item And, solving\\ 		
		$\frac{\partial S}{\partial {\beta_0}} = 
		\sum\nolimits_{i=1}^{n}{\frac{\partial {{\left[ 
		y_i-(\beta_0+\beta_1 x_i) \right]}^{2}}}{\partial {{\beta 
		}_{0}}}}=0 $\\ 
		 
		$ \frac{\partial S}{\partial {{\beta }_{1}}}=\sum\nolimits_{i=1}^{n}{\frac{\partial 
		{{\left[ {{y}_{i}}-({{\beta }_{0}}+{{\beta }_{1}}{{x}_{i}}) 
		\right]}^{2}}}{\partial {{\beta }_{1}}}} = 0$ \\
		
		\vspace{6pt}
		just boils down to solving two simultaneous equations because $ 
		\frac{\partial r_i}{\partial \beta_j} $ is simple {\it because} the model 
		is intrinsically linear:\\
		
		\vspace{6pt}
		$-n{{\beta }_{0}}+\sum\nolimits_{i=1}^{n}{{{y}_{i}}}+{{\beta 
		}_{1}}\sum\nolimits_{i=1}^{n}{{{x}_{i}}}=0 $\\ 
		
		$ \sum\nolimits_{i=1}^{n}{{{x}_{i}}{{y}_{i}}}-{{\beta 
		}_{0}}\sum\nolimits_{i=1}^{n}{{{x}_{i}}}+{{\beta 
		}_{1}}\sum\nolimits_{i=1}^{n}{x_{i}^{2}}=0$

		\item That is, we need to solve $\mathbf{Y} = \mathbf{X 
		\boldsymbol\beta  +  \boldsymbol\varepsilon}$ ({\it this is what R 
		solves when you use {\tt lm()}})

	\end{itemize}

\end{frame}


%%%%%%%%%%%%%%%%%%%%%%%%%%%%%%%%%%%%%%%%%%%%%%%%%%%%%%
\end{document}
